%% ma: L12-B14-0137
%% lop: 12
%% chuong: 5
%% bai: 14
%% dang: 12.14.41
%% loai: TN4
%% muc: NB
%% dap_an: "B"
%% nguon: "Ảnh giáo viên gửi 25/11/2026 (tài liệu ôn tập Phương trình mặt phẳng), Phần I Câu 44"
%% kien_thuc: "Bài 14 (giao điểm của mặt phẳng với trục tọa độ)"
%% trang_thai: chinh-thuc
\begin{ex}
Trong không gian $Oxyz$, cho mặt phẳng $(P)\colon x-2y+3z-6=0$. Giao điểm của mặt phẳng $(P)$ và trục $Ox$ có tọa độ là
\choice{$(0;3;2)$}{\True $(6;0;0)$}{$(2;0;0)$}{$(1;2;-3)$}
\loigiai{Giao điểm với $Ox$ có $y=z=0$. Thay vào $(P)$: $x-6=0\Rightarrow x=6$. Vậy giao điểm là $(6;0;0)$.}
\end{ex}
